% Section 6 - Namespaces
% Roberto Masocco <roberto.masocco@uniroma2.it>
% April 12, 2026

% ### Namespaces ###
\section{Namespaces}
\graphicspath{{figs/section6/}}

% --- Namespaces ---
\begin{frame}{Namespaces}{The naming collision problem}
	\justifying
	C has a single \textbg{global namespace}: every function, variable, and type defined at file scope lives in the same flat space.\\
	\bigskip
	In small programs this is manageable. In large codebases that integrate \textbg{multiple libraries} — written by different authors, at different times — it becomes a \textbg{problem}:
	\begin{itemize}
		\item \textbg{two libraries} may define a function named \texttt{init()}, a struct named \texttt{Node}, or a constant named \texttt{MAX\_SIZE};
		\item the linker cannot resolve the \textbg{collision};
		\item C works around this with \textbg{name prefixes} by convention: \texttt{pthread\_create}, \texttt{PTHREAD\_MUTEX\_INITIALIZER}, \texttt{sched\_setscheduler}...
	\end{itemize}
	\bigskip
	\begin{block}{}
		\centering
		C++ solves this with \textbf{namespaces} — a language-level \textbf{partitioning} of the global name space.
	\end{block}
\end{frame}
\begin{frame}{Namespaces}{Namespaces in C++}
	\justifying
	A \textbg{namespace} is a named scope that groups related declarations:
	\begin{itemize}
		\item defined with the \texttt{namespace} keyword;
		\item may contain functions, classes, variables, type aliases, and other namespaces;
		\item may be \textbg{reopened} across multiple files — all declarations in the same named namespace belong to the same scope.
	\end{itemize}
	\bigskip
	Names inside a namespace are accessed via the \textbg{scope resolution operator} \texttt{::}
	\begin{itemize}
		\item \texttt{std::vector}, \texttt{std::cout}, \texttt{std::string} — the C++ standard library lives in \texttt{std};
		\item \texttt{rclcpp::Node}, \texttt{rclcpp::Publisher} — the ROS 2 client library lives in \texttt{rclcpp}.
	\end{itemize}
	\bigskip
	Namespaces may be \textbg{nested}: \texttt{rclcpp::executors::MultiThreadedExecutor} — each \texttt{::} steps one level deeper, just like dots in a domain name (\emph{e.g.}, \texttt{ing.uniroma2.it}).
\end{frame}
\begin{frame}{Namespaces}{\texttt{using} declarations and directives}
	\justifying
	Writing the full namespace prefix every time may become verbose.\\
	C++ provides two mechanisms to shorten names.\\
	\bigskip
	A \texttt{using} \textbg{declaration introduces a single name} into the current scope:
	\begin{itemize}
		\item \texttt{using std::vector;} — after this, \texttt{vector} alone refers to \texttt{std::vector}.
	\end{itemize}
	\bigskip
	A \texttt{using namespace} \textbg{directive imports all names} from a namespace:
	\begin{itemize}
		\item \texttt{using namespace std;} (\texttt{std} convenient to write but \textbg{dangerous, never do this}).
	\end{itemize}
\end{frame}
\begin{frame}{Namespaces}{Taming long names}
	\justifying
	Namespaces and templates combined may produce \textbg{very long type names}.\\
	C++ provides two ways to create \textbg{aliases}:
	\begin{itemize}
		\item \texttt{typedef} inherited from C: \texttt{typedef std::vector<int> IntVec;}
		\item \texttt{using} alias (preferred in modern C++): \texttt{using IntVec = std::vector<int>;}
	\end{itemize}
	\bigskip
	The \texttt{using} form is preferred because it is \textbg{more readable}, supports templates, and is consistent with the rest of modern C++ syntax.\\
	\bigskip
	You will encounter both in existing code — especially in ROS 2, where aliases like:
	\begin{itemize}
		\item \texttt{using SharedPtr = std::shared\_ptr<Node>;}
	\end{itemize}
	are defined in almost every class to shorten usage at the call site.
\end{frame}
